\begin{proof}[Proof of \cref{obj:thm:normalized-converse}]
Fix \(\beta\in(0,1/4)\). We first construct the required priors for all sufficiently large sample sizes, with a threshold depending only on \(\beta\), and then extend the construction to every positive sample size.

\begin{enumerate}
\item
Use the construction and notation of
\cref{obj:def:paired-prior-handle}. Throughout the paired construction, expectations and probabilities refer to its independent latent draws. The elementary scale bounds are
\[
L\ge1,\qquad 2\le K,\qquad 8L\le K\le9L,\qquad M\le nL.
\]
Indeed, the upper bound for \(K=\lceil8L\rceil\) follows from
\(K<8L+1\le9L\). Consequently,
\[
0\le 2MhB_0=\frac{M}{512nK}\le\frac12,
\qquad
f=1-2MhB_0\ge\frac12,
\qquad
J=f+2\sum_{j=1}^{M}p_j\ge\frac12.
\]
Thus normalization is well defined, and
\[
\frac fJ+2\sum_{j=1}^{M}\frac{p_j}{J}=1,
\qquad
2M+1\le d.
\]
These conclusions also cover \(M=0\), when the law consists of the filler label.

For every \(\sigma,z_j,u\in\{-1,1\}\), the probabilities specified in the construction satisfy
\[
q\le \rho_{\sigma,j,u}\le1,
\qquad
0\le \frac b2+c_0u\le q\le\rho_{\sigma,j,u},
\qquad
0\le\mu_{\sigma,j,u}\le1.
\]
Here the middle bound follows by substituting
\(b=(1+q)/2\), \(c_0=1/16\), and \(q\ge1/2\).
The filler mean is valid as well:
\[
0\le b^2=\frac{(1+q)^2}{4}
\le\frac{1+3q}{4}\le2q,
\qquad
0\le\mu_0=\frac{b^2}{2q}\le1.
\]
Treatment is an independent fair coin, and the realized measurements are selected from their potential measurements as specified in the construction. Conditional on an occupied baseline label and treatment arm, arrival is independent of the outcome: it is complete at the filler label and in the control arm, and has probability \(\rho_{\sigma,j,u}\) in a treated pair cell. Unoccupied cells have zero mass. Hence every constructed normalized law belongs to
\(\mathcal M(d,q)\) in \cref{obj:def:law-class}. The finitely many latent configurations therefore induce discrete probability priors supported on this class.
% lean: CausalSmith.Stat.MarNearcompleteFrontier.pairedFullLaw_legal

\item
We record the interpolation identities needed for the count comparison. Polynomial interpolation at zero gives
\[
\varphi_{\mathrm{pol}}(0)
=\sum_{i=0}^{m}\varphi_{\mathrm{pol}}(x_i)l_i(0)
\qquad
\text{for every polynomial \(\varphi_{\mathrm{pol}}\) of degree at most \(m=K-1\)}.
\]
Applying this identity to the constant polynomial and to positive powers yields
\[
\sum_{i=0}^{m}l_i(0)=1,
\qquad
\sum_{i=0}^{m}\nu_i=\frac1D,
\qquad
\sum_{i=0}^{m}\nu_i x_i^{t_{\mathrm{pol}}}=0
\quad(1\le t_{\mathrm{pol}}\le K-1).
\]
Also,
\[
D=\sum_{i=0}^{m}|l_i(0)|\ge1,
\qquad
\sum_{i=0}^{m}|\nu_i|=1.
\]
Since \(x_i\ge h>0\), the latent probabilities obey
\[
\sum_{i=0}^{m}\frac{h|\nu_i|}{x_i}\le1.
\]
Their residual mass is therefore nonnegative. Substitution into the discrete latent distribution gives
\[
\begin{aligned}
\mathbb E p&=hB_0,
&
\mathbb E(pz)&=\frac{hB_0}{D},\\
\mathbb E(p^{t_{\mathrm{pol}}}z)
&=hB_0^{t_{\mathrm{pol}}}
  \sum_{i=0}^{m}\nu_i x_i^{t_{\mathrm{pol}}-1}
 =0
&& (2\le t_{\mathrm{pol}}\le K),\\
\mathbb E p^2
&=\mathbb E[(pz)^2]
 =hB_0^2\sum_{i=0}^{m}|\nu_i|x_i
 \le hB_0^2
 =\frac1{1024^2n^2}.
\end{aligned}
\]
The atom with \(p=0\) contributes zero to each displayed positive-power moment.
% lean: CausalSmith.Stat.MarNearcompleteFrontier.latent_signed_power_moment_zero

\item
Let the eight-coordinate pair vector be
\[
V_\sigma(z,r)
=\bigl(v_\sigma(z,r),v_\sigma(z,-r)\bigr),
\qquad
\sigma,z,r\in\{-1,1\}.
\]
Its coordinates are nonnegative and satisfy
\[
\sum_{\ell_{\mathrm{atom}}=1}^{8}
V_{\sigma,\ell_{\mathrm{atom}}}(z,r)=2.
\]
For a count vector
\(\boldsymbol{k}_{\mathrm{pair}}
=(k_{\mathrm{pair},1},\ldots,k_{\mathrm{pair},8})\in\mathbb N^8\),
define
\[
t_{\mathrm{count}}
=\sum_{\ell_{\mathrm{atom}}=1}^{8}k_{\mathrm{pair},\ell_{\mathrm{atom}}},
\qquad
\mathcal H_\sigma(z;\boldsymbol{k}_{\mathrm{pair}})
=\frac12\sum_{r\in\{-1,1\}}
 \prod_{\ell_{\mathrm{atom}}=1}^{8}
 V_{\sigma,\ell_{\mathrm{atom}}}(z,r)^{
 k_{\mathrm{pair},\ell_{\mathrm{atom}}}}
\]
Products use the convention \(0^0=1\). The likelihood of this count vector in one mixed pair block is
\[
\mathcal L_{\sigma,\mathrm{pair}}
(\boldsymbol{k}_{\mathrm{pair}})
=
\mathbb E\!\left[
e^{-8np}\,
\frac{(4np)^{t_{\mathrm{count}}}}
     {\prod_{\ell_{\mathrm{atom}}=1}^{8}
       (k_{\mathrm{pair},\ell_{\mathrm{atom}}})!}
\,
\mathcal H_\sigma(z;\boldsymbol{k}_{\mathrm{pair}})
\right],
\]
where the expectation is over the latent pair \((p,z)\).

Changing \(z\) to \(-z\) exchanges the two signs, so
\[
\mathcal H_-(z;\boldsymbol{k}_{\mathrm{pair}})
-\mathcal H_+(z;\boldsymbol{k}_{\mathrm{pair}})
=
z\bigl[
\mathcal H_-(1;\boldsymbol{k}_{\mathrm{pair}})
-\mathcal H_+(1;\boldsymbol{k}_{\mathrm{pair}})
\bigr].
\]
This difference is zero when \(t_{\mathrm{count}}=0\). It is also zero when
\(t_{\mathrm{count}}=1\), because averaging over the fair orientation gives, at either label, the four coefficients
\[
\left(\frac12,\frac\delta4,\frac b4,\frac b4\right),
\]
independently of \(\sigma z\).

For \(v_{\mathrm{exp}}\in\mathbb N\), define the signed Taylor coefficient
\[
\Gamma_{v_{\mathrm{exp}}}
(\boldsymbol{k}_{\mathrm{pair}})
=
\mathbb E\!\left[
\frac{(-8np)^{v_{\mathrm{exp}}}}{v_{\mathrm{exp}}!}
\frac{(4np)^{t_{\mathrm{count}}}}
     {\prod_{\ell_{\mathrm{atom}}=1}^{8}
       (k_{\mathrm{pair},\ell_{\mathrm{atom}}})!}
\bigl(\mathcal H_-(z;\boldsymbol{k}_{\mathrm{pair}})
-\mathcal H_+(z;\boldsymbol{k}_{\mathrm{pair}})\bigr)
\right].
\]
The exponential series, together with \(0\le p\le B_0\), gives
\[
\mathcal L_{-,\mathrm{pair}}(\boldsymbol{k}_{\mathrm{pair}})
-\mathcal L_{+,\mathrm{pair}}(\boldsymbol{k}_{\mathrm{pair}})
=
\sum_{v_{\mathrm{exp}}=0}^{\infty}
\Gamma_{v_{\mathrm{exp}}}(\boldsymbol{k}_{\mathrm{pair}}).
\]
The low-count cancellations above and the signed moment identities imply
\[
\Gamma_{v_{\mathrm{exp}}}(\boldsymbol{k}_{\mathrm{pair}})=0
\qquad
\text{if }v_{\mathrm{exp}}+t_{\mathrm{count}}\le K.
\]
Indeed, when \(t_{\mathrm{count}}\ge2\), this coefficient is a fixed multiple of
\(\mathbb E[p^{v_{\mathrm{exp}}+t_{\mathrm{count}}}z]\).

For fixed count degree, the multinomial identity gives
\[
\sum_{\boldsymbol{k}_{\mathrm{pair}}:\,
      t_{\mathrm{count}}=t_{\mathrm{degree}}}
\frac{\mathcal H_\sigma(z;\boldsymbol{k}_{\mathrm{pair}})}
     {\prod_{\ell_{\mathrm{atom}}=1}^{8}
       (k_{\mathrm{pair},\ell_{\mathrm{atom}}})!}
=\frac{2^{t_{\mathrm{degree}}}}{t_{\mathrm{degree}}!},
\qquad t_{\mathrm{degree}}\in\mathbb N.
\]
Bounding the absolute mark difference by the sum of its two nonnegative terms therefore yields
\[
\sum_{\boldsymbol{k}_{\mathrm{pair}}:\,
      t_{\mathrm{count}}=t_{\mathrm{degree}}}
|\Gamma_{v_{\mathrm{exp}}}(\boldsymbol{k}_{\mathrm{pair}})|
\le
2\frac{(8nB_0)^{v_{\mathrm{exp}}}}{v_{\mathrm{exp}}!}
 \frac{(8nB_0)^{t_{\mathrm{degree}}}}{t_{\mathrm{degree}}!}.
\]
Thus, for each total degree \(s_{\mathrm{tot}}\in\mathbb N\),
\[
\sum_{t_{\mathrm{degree}}=0}^{s_{\mathrm{tot}}}
\ \sum_{\boldsymbol{k}_{\mathrm{pair}}:\,
      t_{\mathrm{count}}=t_{\mathrm{degree}}}
|\Gamma_{s_{\mathrm{tot}}-t_{\mathrm{degree}}}(\boldsymbol{k}_{\mathrm{pair}})|
\le
\begin{cases}
0,&s_{\mathrm{tot}}\le K,\\[2mm]
\displaystyle
2\frac{(16nB_0)^{s_{\mathrm{tot}}}}{s_{\mathrm{tot}}!},
&s_{\mathrm{tot}}>K.
\end{cases}
\]
The second line follows by factorial-series convolution. Its summable majorant justifies the absolute summations and gives
\[
\sum_{\boldsymbol{k}_{\mathrm{pair}}\in\mathbb N^8}
\left|
\mathcal L_{-,\mathrm{pair}}(\boldsymbol{k}_{\mathrm{pair}})
-\mathcal L_{+,\mathrm{pair}}(\boldsymbol{k}_{\mathrm{pair}})
\right|
\le
2\sum_{s_{\mathrm{tot}}>K}
\frac{(16nB_0)^{s_{\mathrm{tot}}}}{s_{\mathrm{tot}}!}.
\]

Denote the two complete mixed count laws by
\[
\mathsf C_\sigma
=\operatorname{Law}(\text{the unnormalized observed-atom counts
under sign }\sigma),
\qquad \sigma\in\{-1,1\}.
\]
The independent pair blocks have probability likelihoods, and the filler likelihood is common to both signs. Write the tuple of pair count vectors as
\[
\boldsymbol{k}_{\mathrm{blocks}}
=\bigl(\boldsymbol{k}_{\mathrm{pair}}^{(1)},\ldots,
       \boldsymbol{k}_{\mathrm{pair}}^{(M)}\bigr)
\in(\mathbb N^8)^M.
\]
Each block likelihood is nonnegative and normalized:
\[
\sum_{\boldsymbol{k}_{\mathrm{pair}}\in\mathbb N^8}
\mathcal L_{\sigma,\mathrm{pair}}(\boldsymbol{k}_{\mathrm{pair}})=1,
\qquad \sigma\in\{-1,1\}.
\]
The product difference has the telescoping identity
\[
\begin{aligned}
&\prod_{j=1}^{M}
\mathcal L_{-,\mathrm{pair}}(\boldsymbol{k}_{\mathrm{pair}}^{(j)})
-\prod_{j=1}^{M}
\mathcal L_{+,\mathrm{pair}}(\boldsymbol{k}_{\mathrm{pair}}^{(j)})\\
&=\sum_{j=1}^{M}
\Bigl[
\mathcal L_{-,\mathrm{pair}}(\boldsymbol{k}_{\mathrm{pair}}^{(j)})
-\mathcal L_{+,\mathrm{pair}}(\boldsymbol{k}_{\mathrm{pair}}^{(j)})
\Bigr]
\prod_{\ell_{\mathrm{block}}=1}^{j-1}
\mathcal L_{-,\mathrm{pair}}
 (\boldsymbol{k}_{\mathrm{pair}}^{(\ell_{\mathrm{block}})})
\prod_{\ell_{\mathrm{block}}=j+1}^{M}
\mathcal L_{+,\mathrm{pair}}
 (\boldsymbol{k}_{\mathrm{pair}}^{(\ell_{\mathrm{block}})}).
\end{aligned}
\]
Empty products equal one. Taking absolute values and summing over all block count vectors gives
\[
\begin{aligned}
&\sum_{\boldsymbol{k}_{\mathrm{blocks}}\in(\mathbb N^8)^M}
\left|
\prod_{j=1}^{M}
\mathcal L_{-,\mathrm{pair}}(\boldsymbol{k}_{\mathrm{pair}}^{(j)})
-\prod_{j=1}^{M}
\mathcal L_{+,\mathrm{pair}}(\boldsymbol{k}_{\mathrm{pair}}^{(j)})
\right|\\
&\le\sum_{j=1}^{M}
\sum_{\boldsymbol{k}_{\mathrm{pair}}\in\mathbb N^8}
\left|
\mathcal L_{-,\mathrm{pair}}(\boldsymbol{k}_{\mathrm{pair}})
-\mathcal L_{+,\mathrm{pair}}(\boldsymbol{k}_{\mathrm{pair}})
\right|\\
&\le 2M\sum_{s_{\mathrm{tot}}>K}
\frac{(16nB_0)^{s_{\mathrm{tot}}}}{s_{\mathrm{tot}}!}.
\end{aligned}
\]
Indeed, in each telescoping summand, summing every block other than the selected block gives one by normalization. The sum over the \(M\) selected blocks therefore gives the factor \(M\). The common filler likelihood also sums to one. Splitting a full count vector into these blocks consequently gives
\[
\operatorname{TV}(\mathsf C_-,\mathsf C_+)
\le
\frac12\sum_{\text{full count vectors}}
|\mathsf C_-(\{\text{vector}\})-\mathsf C_+(\{\text{vector}\})|
\le
M\sum_{s_{\mathrm{tot}}>K}
\frac{(16nB_0)^{s_{\mathrm{tot}}}}{s_{\mathrm{tot}}!}.
\]
Counts outside the specified atoms are zero under both laws. For \(M=0\), both count laws equal the common filler law and the bound is zero.
% lean: CausalSmith.Stat.MarNearcompleteFrontier.mixedCountLaw_tv_le_factorial_tail

\item
For \(\xi\ge0\), define
\[
\mathcal E_K(\xi)
=\sum_{s_{\mathrm{tot}}>K}
\frac{\xi^{s_{\mathrm{tot}}}}{s_{\mathrm{tot}}!}.
\]
Termwise exponential weighting gives
\[
\mathcal E_K(\xi)
\le e^{-K}\sum_{s_{\mathrm{tot}}=0}^{\infty}
\frac{(e\xi)^{s_{\mathrm{tot}}}}{s_{\mathrm{tot}}!}
=\exp(-K+e\xi),
\]
because \(e^{s_{\mathrm{tot}}-K}\ge1\) on the tail. Since
\(16nB_0=K/64\), \(e\le3\), and \(K\ge8L\),
\[
\mathcal E_K(16nB_0)
\le\exp(-K+eK/64)\le e^{-4L}.
\]
Using \(M\le nL\), \(n\le e^L\), and \(L\le e^L\), we obtain
\[
\operatorname{TV}(\mathsf C_-,\mathsf C_+)
\le Me^{-4L}
\le nLe^{-4L}
\le e^{-2L}
\le e^{-L}
=\frac1{e+n}
\le\frac1n.
\]
Choose an integer threshold satisfying
\[
N_{\mathrm{count}}\ge2,
\qquad
N_{\mathrm{count}}>\frac2\beta.
\]
Then, uniformly in \(d,q\),
\[
n\ge N_{\mathrm{count}}
\quad\Longrightarrow\quad
\operatorname{TV}(\mathsf C_-,\mathsf C_+)\le\frac\beta2.
\]
% lean: CausalSmith.Stat.MarNearcompleteFrontier.mixedCountLaw_tv_eventual

\item
We next transfer the count comparison to the complete fixed-size sample. Conditional on the latent variables, define the total count and its mean by
\[
N_{\mathrm{Pois}}
=\sum_{o\in O}\text{count at }o,
\qquad
\Lambda_{\mathrm{count}}=4nJ.
\]
The exact observed-atom probabilities identify the count experiment as independent Poisson counts from \(P_{\sigma,O}\) at total intensity
\(\Lambda_{\mathrm{count}}\). For example,
\[
4np_jv_\sigma(z_j,u)
=(4nJ)\frac{p_j}{J}v_\sigma(z_j,u),
\]
and the same identity holds for the filler coefficients. Hence
\[
N_{\mathrm{Pois}}\mid\text{latent variables}
\sim\operatorname{Poisson}(\Lambda_{\mathrm{count}}),
\qquad
\Lambda_{\mathrm{count}}\ge2n.
\]

For this conditional Poisson law, put
\[
x_{\mathrm{tail}}
=\sqrt{\frac{\Lambda_{\mathrm{count}}n}{4}},
\qquad
\vartheta_{\mathrm{tail}}
=\frac{x_{\mathrm{tail}}}{\Lambda_{\mathrm{count}}}.
\]
Both are positive, and
\[
x_{\mathrm{tail}}\le\frac{\Lambda_{\mathrm{count}}}{2}
\le\Lambda_{\mathrm{count}}-n.
\]
The centered Poisson exponential moment is
\[
\mathbb E\!\left[
e^{-\vartheta_{\mathrm{tail}}
(N_{\mathrm{Pois}}-\Lambda_{\mathrm{count}})}
\,\middle|\,\text{latent variables}\right]
=
\exp\!\left\{
\Lambda_{\mathrm{count}}
(e^{-\vartheta_{\mathrm{tail}}}-1+\vartheta_{\mathrm{tail}})
\right\}.
\]
For every nonnegative real \(\vartheta_{\mathrm{tail}}\),
\[
e^{-\vartheta_{\mathrm{tail}}}-1+\vartheta_{\mathrm{tail}}
\le\frac{\vartheta_{\mathrm{tail}}^2}{2}.
\]
One way to verify this bound is to combine
\(e^{\vartheta_{\mathrm{tail}}}\ge
1+\vartheta_{\mathrm{tail}}+\vartheta_{\mathrm{tail}}^2/2\)
with
\[
\left(1+\vartheta_{\mathrm{tail}}
+\frac{\vartheta_{\mathrm{tail}}^2}{2}\right)
\left(1-\vartheta_{\mathrm{tail}}
+\frac{\vartheta_{\mathrm{tail}}^2}{2}\right)
=1+\frac{\vartheta_{\mathrm{tail}}^4}{4}\ge1.
\]
The second factor is positive, so taking reciprocals gives the claimed bound. Exponential Markov inequality now yields
\[
\begin{aligned}
\Pr(N_{\mathrm{Pois}}<n\mid\text{latent variables})
&\le
\Pr(\Lambda_{\mathrm{count}}-N_{\mathrm{Pois}}>
x_{\mathrm{tail}}\mid\text{latent variables})\\
&\le
\exp\!\left\{
-\frac{x_{\mathrm{tail}}^2}{\Lambda_{\mathrm{count}}}
+\Lambda_{\mathrm{count}}
(e^{-\vartheta_{\mathrm{tail}}}-1+\vartheta_{\mathrm{tail}})
\right\}\\
&\le
\exp\!\left(-\frac{x_{\mathrm{tail}}^2}
                   {2\Lambda_{\mathrm{count}}}\right)
=e^{-n/8}
\le\frac8n.
\end{aligned}
\]

Uniformly ordering independent Poisson counts produces iid marks with law
\(P_{\sigma,O}\), independently of their total count. Consequently,
conditional on the latent variables, the output of
\(\kappa_n\) on \(N_{\mathrm{Pois}}\ge n\) has law
\(P_{\sigma,O}^{\otimes n}\). Coupling this output with an iid sample and using the fixed fallback on the complementary event gives
\[
\operatorname{TV}\bigl(Q_\sigma^{(n)},\kappa_n\mathsf C_\sigma\bigr)
\le
\mathbb E\bigl[
\Pr(N_{\mathrm{Pois}}<n\mid\text{latent variables})\bigr]
\le\frac8n.
\]
Here \(\kappa_n\mathsf C_\sigma\) denotes the output law of the kernel.
Its independence of the sign and latent variables permits total-variation contraction. The triangle inequality therefore gives
\[
\operatorname{TV}(Q_-^{(n)},Q_+^{(n)})
\le
\operatorname{TV}(\mathsf C_-,\mathsf C_+)+\frac{16}{n}.
\]
Choose integer thresholds
\[
N_{\mathrm{transfer}}\ge2,
\qquad
N_{\mathrm{transfer}}>\frac{32}{\beta},
\qquad
N_{\mathrm{TV}}
=\max\{N_{\mathrm{count}},N_{\mathrm{transfer}}\}.
\]
For every \(n\ge N_{\mathrm{TV}}\), the two terms on the right are each at most \(\beta/2\), so
\[
\operatorname{TV}(Q_-^{(n)},Q_+^{(n)})\le\beta.
\]
% lean: CausalSmith.Stat.MarNearcompleteFrontier.paired_predictive_tv_eventual

\item
We now establish target separation. First, the interpolation norm has the uniform bound
\[
1\le D\le\cosh2.
\]
To see the upper bound, define the transformed nodes, exterior point, and hyperbolic parameter by
\[
y_i=\frac{2x_i-1-h}{1-h}
=\cos\frac{i\pi}{m},
\qquad
y_*=-\frac{1+h}{1-h},
\qquad
t_{\mathrm{hyp}}=\log\frac{K+1}{K-1}.
\]
Use the Chebyshev convention
\[
T_0(t_{\mathrm{Cheb}})=1,\qquad
T_1(t_{\mathrm{Cheb}})=t_{\mathrm{Cheb}},\qquad
T_{\ell_{\mathrm{pol}}+1}(t_{\mathrm{Cheb}})
=2t_{\mathrm{Cheb}}T_{\ell_{\mathrm{pol}}}(t_{\mathrm{Cheb}})
-T_{\ell_{\mathrm{pol}}-1}(t_{\mathrm{Cheb}})
\quad(\ell_{\mathrm{pol}}\ge1,\ t_{\mathrm{Cheb}}\in\mathbb R).
\]
The ordered nodes give
\(\operatorname{sign}(l_i(0))=(-1)^{m+i}\), while
\(T_m(y_i)=(-1)^i\). Interpolating \(T_m\) at \(y_*\) therefore gives
\[
D=|T_m(y_*)|.
\]
Since \(h=K^{-2}\),
\[
y_*=-\cosh(t_{\mathrm{hyp}}),
\qquad
0\le mt_{\mathrm{hyp}}
=(K-1)\log\frac{K+1}{K-1}
\le
(K-1)\left(\frac{K+1}{K-1}-1\right)=2.
\]
The inequality uses \(\log t\le t-1\) for \(t>0\).
The Chebyshev hyperbolic identity then gives
\[
D=\cosh(mt_{\mathrm{hyp}})\le\cosh2.
\]

Define the signed latent mass, its mean, and the fixed hyperbolic constant by
\[
W=\sum_{j=1}^{M}p_jz_j,
\qquad
\mu_W=\frac{MhB_0}{D}
=\frac{M}{1024nKD},
\qquad
C_{\mathrm{hyp}}=\cosh2.
\]
Independence of the latent pairs and the moment identities above imply
\[
\mathbb EW=\mu_W,
\qquad
\operatorname{Var}(W)
=M\left(\mathbb E[(pz)^2]-\left(\frac{hB_0}{D}\right)^2\right)
\le\frac{M}{1024^2n^2}.
\]
Similarly, the deterministic filler mass gives
\[
J=1+2\sum_{j=1}^{M}(p_j-hB_0),
\qquad
\mathbb EJ=1,
\]
and
\[
\operatorname{Var}(J)
=4M\bigl(\mathbb Ep^2-(hB_0)^2\bigr)
\le\frac{4M}{1024^2n^2}.
\]
In both variance computations, the cross terms vanish because independent centered summands have zero product expectation.

Finally, opposite orientations in a pair satisfy the exact identity
\[
\mu_{\sigma,j,1}+\mu_{\sigma,j,-1}
=2\mu_0+\frac{\sigma\delta z_j}{16q}.
\]
This follows by putting the two fractions over a common denominator and using
\[
b^2-\frac{\delta^2}{4}=q.
\]
Since \(Y(0)=0\), summing the filler and pair contributions and dividing by \(J\) yields
\[
\begin{aligned}
\tau(P_\sigma)
&=
\frac{
f\mu_0+\sum_{j=1}^{M}
p_j\left(2\mu_0+\frac{\sigma\delta z_j}{16q}\right)}
{J}\\
&=\mu_0+\frac{\sigma\delta}{16q}\frac WJ.
\end{aligned}
\]
% lean: CausalSmith.Stat.MarNearcompleteFrontier.pairedFullLaw_tau_exact

\item
For \(n,d\ge1\), define the dimension factor
\[
\zeta_{n,d}
=\min\left\{1,\frac d{nL}\right\}\in[0,1],
\qquad
g_{n,d,q}=\delta\zeta_{n,d}.
\]
Under the stated separation regime,
\[
\frac1n\le\delta^2\zeta_{n,d}^2
\le\frac14\zeta_{n,d}^2,
\qquad
\frac4n\le\zeta_{n,d}^2.
\]
In particular, \(\zeta_{n,d}>0\). For \(n\ge2\), the regime also implies \(d\ge3\): if \(d\le2\), then
\[
0\le g_{n,d,q}
\le\frac12\frac2{nL}\le\frac1n,
\qquad
g_{n,d,q}^2\le\frac1{n^2}<\frac1n,
\]
a contradiction.

For \(d\ge3\) and \(nL\ge2\), the integer floors satisfy
\[
\left\lfloor\frac{d-1}{2}\right\rfloor\ge\frac d4,
\qquad
\lfloor nL\rfloor\ge\frac{nL}{2}.
\]
Thus
\[
M\ge\frac14\min\{d,nL\}
=\frac{nL}{4}\zeta_{n,d}>0,
\]
and, using \(K\le9L\),
\[
\frac{M}{nK}\ge\frac1{36}\zeta_{n,d},
\qquad
\mu_W\ge
\frac1{36\cdot1024\,C_{\mathrm{hyp}}}\zeta_{n,d}>0.
\]
Squaring the lower bound for \(M\) and using
\(\zeta_{n,d}^2\ge4/n\) also gives
\[
M^2\ge\frac{nL^2}{4},
\qquad
K^4\le6561L^4,
\qquad
K^4\le\frac{26244L^2}{n}M^2.
\]

The logarithmic factor obeys
\[
\forall\varepsilon>0\ \exists N_\varepsilon\in\mathbb N\
\forall n\ge N_\varepsilon:
\qquad L^2\le\varepsilon n.
\]
Indeed, put \(\gamma_{\log}=e+1\). The standard limit
\((\log \xi)^2/\xi\to0\) as \(\xi\to\infty\) gives, eventually,
\[
L^2\le\frac{\varepsilon}{\gamma_{\log}}(e+n)
\le\varepsilon n,
\]
where the last inequality uses \(n\ge1\).

Define
\[
\eta_W=\frac{\beta}{8C_{\mathrm{hyp}}^2}>0.
\]
Choose \(N_W\ge2\), depending only on \(\beta\), so that
\[
n\ge N_W
\quad\Longrightarrow\quad
L^2\le\frac{\eta_W^2}{26244}n.
\]
For every admissible \(n,d,q\) beyond this threshold,
\[
K^4\le\eta_W^2M^2,
\qquad
\frac{K^2}{M}\le\eta_W.
\]
The second inequality follows by taking nonnegative square roots.
% lean: CausalSmith.Stat.MarNearcompleteFrontier.eventually_priorK_sq_div_pairCount_le

\item
Chebyshev's inequality and the positive mean established above give, for admissible \(n\ge N_W\),
\[
\begin{aligned}
\Pr\left(|W-\mu_W|>\frac{\mu_W}{2}\right)
&\le
\frac{M/(1024^2n^2)}{(\mu_W/2)^2}\\
&=\frac{4K^2D^2}{M}
\le4C_{\mathrm{hyp}}^2\frac{K^2}{M}
\le\frac\beta2.
\end{aligned}
\]
Choose another integer threshold \(N_J\), depending only on \(\beta\), such that
\[
n\ge N_J
\quad\Longrightarrow\quad
L^2\le\frac{\beta\,1024^2}{32}n.
\]
The normalizer variance bound, \(M\le nL\), and \(L\ge1\) then yield
\[
\begin{aligned}
\Pr\left(|J-1|>\frac12\right)
&\le\frac{16M}{1024^2n^2}\\
&\le\frac{16L}{1024^2n}
\le\frac{16L^2}{1024^2n}
\le\frac\beta2.
\end{aligned}
\]

Define the common good event, target threshold, and separation constant by
\[
\mathcal G
=\left\{|W-\mu_W|\le\frac{\mu_W}{2},
          \ |J-1|\le\frac12\right\},
\qquad
N_{\mathrm{target}}=\max\{N_W,N_J\},
\]
\[
c_{\mathrm{large}}
=\frac1{36\cdot1024\,C_{\mathrm{hyp}}\cdot48}>0.
\]
For admissible \(n\ge N_{\mathrm{target}}\), the union bound gives
\[
\Pr(\mathcal G)\ge1-\beta.
\]
On \(\mathcal G\),
\[
W\ge\frac{\mu_W}{2},
\qquad
0<J\le\frac32,
\qquad
\frac WJ\ge\frac{\mu_W}{3}.
\]
Since \(\delta\ge0\) and \(0<q\le1\), it follows that
\[
\frac{\delta}{16q}\frac WJ
\ge\frac{\delta\mu_W}{48}
\ge
\frac{\delta\zeta_{n,d}}
     {36\cdot1024\,C_{\mathrm{hyp}}\cdot48}
=c_{\mathrm{large}}g_{n,d,q}.
\]
Set
\[
m_0=\mu_0.
\]
The exact target identity therefore implies, on \(\mathcal G\),
\[
\tau(P_-)\le m_0-c_{\mathrm{large}}g_{n,d,q},
\qquad
m_0+c_{\mathrm{large}}g_{n,d,q}\le\tau(P_+).
\]
Because \(\Pi_\sigma=\operatorname{Law}(P_\sigma)\), the corresponding prior event masses equal the probabilities of their latent preimages. Hence
\[
\Pi_-\{\tau(P)\le m_0-c_{\mathrm{large}}g_{n,d,q}\}
\ge1-\beta,
\]
\[
\Pi_+\{m_0+c_{\mathrm{large}}g_{n,d,q}\le\tau(P)\}
\ge1-\beta.
\]
% lean: CausalSmith.Stat.MarNearcompleteFrontier.pairedPrior_target_events_eventual

\item
Combine the two large-sample thresholds:
\[
N_\beta=\max\{N_{\mathrm{TV}},N_{\mathrm{target}}\}\ge2.
\]
For every admissible \(n\ge N_\beta\), the paired priors simultaneously satisfy
\[
\operatorname{TV}(Q_-^{(n)},Q_+^{(n)})\le\beta
\]
and both target-event inequalities with separation constant
\(c_{\mathrm{large}}\). The threshold and this positive constant are independent of \(n,d,q\).
% lean: CausalSmith.Stat.MarNearcompleteFrontier.normalized_converse

\item
To cover smaller sample sizes, define
\[
a_\beta
=\min\left\{\frac14,
\frac{\sqrt{\log(1+4\beta^2)}}4\right\}>0,
\qquad
u_n=\frac{a_\beta}{\sqrt n},
\qquad n\ge1.
\]
These choices satisfy
\[
0<u_n\le\frac14,
\qquad
16nu_n^2=16a_\beta^2\le\log(1+4\beta^2).
\]

Choose one baseline label, set both potential surrogates and \(Y(0)\) to zero, assign treatment by an independent fair coin, and make arrival complete in both arms. Let
\(P_{-,\mathrm{small}}\) and \(P_{+,\mathrm{small}}\) be the resulting laws with
\[
Y(1)\sim\operatorname{Bernoulli}(1/2)
\quad\text{and}\quad
Y(1)\sim\operatorname{Bernoulli}(1/2+u_n),
\]
respectively. Both laws belong to \(\mathcal M(d,q)\) for every \(d\ge1\) and \(q\in[1/2,1]\), and
\[
\tau(P_{-,\mathrm{small}})=\frac12,
\qquad
\tau(P_{+,\mathrm{small}})=\frac12+u_n.
\]
In the observed-atom order
\((A,R,RY)=(0,1,0),(1,1,0),(1,1,1)\), their probabilities are
\[
P_{-,\mathrm{small},O}:
\left(\frac12,\frac14,\frac14\right),
\qquad
P_{+,\mathrm{small},O}:
\left(\frac12,\frac14-\frac{u_n}{2},
                 \frac14+\frac{u_n}{2}\right).
\]
All other atoms have zero mass. Thus the plus observed law is absolutely continuous with respect to the minus observed law, and direct summation gives
\[
\begin{aligned}
\chi^2(P_{+,\mathrm{small},O}\Vert P_{-,\mathrm{small},O})
&:=
\sum_{\substack{o\in O\\
P_{-,\mathrm{small},O}(\{o\})>0}}
\frac{
\bigl(P_{+,\mathrm{small},O}(\{o\})
-P_{-,\mathrm{small},O}(\{o\})\bigr)^2}
{P_{-,\mathrm{small},O}(\{o\})}\\
&=2u_n^2\le16u_n^2.
\end{aligned}
\]
For finite probability laws, multiplication of likelihood ratios gives the product identity for \(1+\chi^2\), while Cauchy--Schwarz gives
\(\operatorname{TV}\le\frac12\sqrt{\chi^2}\).
Applying these facts to the displayed pair yields
\[
\begin{aligned}
\operatorname{TV}
\left(P_{-,\mathrm{small},O}^{\otimes n},
      P_{+,\mathrm{small},O}^{\otimes n}\right)
&\le
\frac12\sqrt{
\left(1+
\chi^2(P_{+,\mathrm{small},O}\Vert
       P_{-,\mathrm{small},O})\right)^n-1}\\
&\le\frac12\sqrt{e^{16nu_n^2}-1}
\le\beta.
\end{aligned}
\]
Here the middle inequality follows from \(1+t\le e^t\) for \(t\ge0\).

Define
\[
c_{\mathrm{small}}=\frac{a_\beta}{\sqrt{N_\beta}}>0.
\]
For \(1\le n<N_\beta\), the scale bound
\(0\le g_{n,d,q}\le1/2\) gives
\[
c_{\mathrm{small}}g_{n,d,q}
\le\frac{a_\beta}{2\sqrt{N_\beta}}
\le\frac{a_\beta}{2\sqrt n}
=\frac{u_n}{2}.
\]
Take the point priors and midpoint
\[
\Pi_-^{\mathrm{small}}=\delta_{P_{-,\mathrm{small}}},
\qquad
\Pi_+^{\mathrm{small}}=\delta_{P_{+,\mathrm{small}}},
\qquad
m_0=\frac12+\frac{u_n}{2}.
\]
Their target inequalities hold deterministically:
\[
\tau(P_{-,\mathrm{small}})
\le m_0-c_{\mathrm{small}}g_{n,d,q},
\qquad
m_0+c_{\mathrm{small}}g_{n,d,q}
\le\tau(P_{+,\mathrm{small}}).
\]
Both prior event masses are therefore one, and the preceding product bound supplies the required mixture distance.
% lean: CausalSmith.Stat.MarNearcompleteFrontier.normalized_converse_smallN

\item
Finally, set
\[
c_\beta=\min\{c_{\mathrm{large}},c_{\mathrm{small}}\}>0.
\]
If \(n\ge N_\beta\), select the paired priors and their center. If
\(1\le n<N_\beta\), select the point priors and their midpoint.
In either case, write
\[
c_{\mathrm{branch}}
=
\begin{cases}
c_{\mathrm{large}},&n\ge N_\beta,\\
c_{\mathrm{small}},&1\le n<N_\beta.
\end{cases}
\]
Since \(g_{n,d,q}\ge0\) and \(c_\beta\le c_{\mathrm{branch}}\),
\[
\{P:\tau(P)\le m_0-c_{\mathrm{branch}}g_{n,d,q}\}
\subseteq
\{P:\tau(P)\le m_0-c_\beta g_{n,d,q}\},
\]
\[
\{P:m_0+c_{\mathrm{branch}}g_{n,d,q}\le\tau(P)\}
\subseteq
\{P:m_0+c_\beta g_{n,d,q}\le\tau(P)\}.
\]
Monotonicity of prior probabilities preserves both lower bounds
\(1-\beta\), and the mixture-distance bounds remain valid. These two cases cover every admissible positive sample size and establish all the assertions.
% lean: CausalSmith.Stat.MarNearcompleteFrontier.normalized_converse_of_eventual_certificate
\end{enumerate}
\end{proof}
